\devoirDS{2}{Chapitres 4 à 6}{Un ion à reconstituer/6,Des électrons aux molécules/7,Compter du dioxyde de carbone/7}
\begin{blocLivre}{ds-2-1}\begin{exoBac}{6}{Un ion à reconstituer}
Un noyau contient 13 protons et 14 neutrons. L’entité possède 10 électrons. L’élément de numéro 13 est l’aluminium ; $e=\qty{1,60e-19}{C}$.
\begin{questions}
\item Écrire le symbole du noyau.
\item Calculer la charge de l’entité et donner sa formule.
\item Écrire le composé neutre avec des ions chlorure.
\end{questions}
\end{exoBac}\end{blocLivre}
\begin{blocLivre}{ds-2-2}\begin{exoBac}{7}{Des électrons aux molécules}
L’oxygène a $Z=8$ ; on fournit les schémas de Lewis.
\schema{lewis}
\begin{questions}
\item Écrire sa configuration électronique et donner ses électrons de valence.
\item Interpréter les doublets autour de O dans l’eau.
\item Quelle configuration de gaz noble sert de comparaison ?
\item Pourquoi rompre les liaisons de la molécule nécessite-t-il un apport d’énergie ?
\end{questions}
\end{exoBac}\end{blocLivre}
\begin{blocLivre}{ds-2-3}\begin{exoBac}{7}{Compter du dioxyde de carbone}
Un échantillon pur contient 0,731 g de \ce{CO2}. $m_C=1,99\times10^{-26}$ kg ; $m_O=2,66\times10^{-26}$ kg ; $N_A=6,02\times10^{23}\unit{mol^{-1}}$.
\begin{questions}
\item Calculer la masse d’une molécule.
\item Calculer N puis n.
\item Donner le nombre d’atomes d’oxygène.
\item Quelle erreur ferait-on en utilisant 0,731 sans conversion dans le quotient par la masse moléculaire ?
\end{questions}
\end{exoBac}\end{blocLivre}

